\documentclass[a4paper,12pt]{article}
\usepackage{docsTemplate}
\usepackage{longtable}
\usepackage{tabularx}
\usepackage[table]{xcolor}

\setTitle{Verbale Interno 23/12/2025}
\setAuthors{Andrea Masiero}
\setVerificators{Bianca Zaghetto}
\setApprovation{Sara Gioia Fichera}
\setVersion{1.0}
\setType{Interno}
\setPlace{Discord}
\setTime{16:00}
\setDestination{Team Three Technologies}
\setAbsents{Nessuno}

\begin{document}
\include{memoTitlepage}

\newpage
\section*{Registro delle modifiche}
\rowcolors{2}{lightgray}{white}
\begin{tabularx}{\textwidth}{|l|l|l|l|X|X|}
\hline
\textbf{Vers.} & \textbf{Data} & \textbf{Autore} & \textbf{Ruolo} & \textbf{Verifica} & \textbf{Oggetto} \\
\hline
1.0 & 28/12/25 & Sara Gioia Fichera & Responsabile &  & Approvazione \\
\hline
0.3 & 28/12/25 & Andrea Masiero & Analista & Mattia Oliva Medin & Correzioni \\
\hline
0.2 & 27/12/25 & Andrea Masiero & Analista & Bianca Zaghetto & Correzioni \\
\hline
0.1 & 27/12/25 & Andrea Masiero & Analista & & Prima stesura \\
\hline
\end{tabularx}

\newpage
\tableofcontents

\newpage
\markright{}

\newpage
\section{Ordine del giorno}
\begin{itemize}
    \item Stato avanzamento Analisi dei Requisiti
    \item Discussione sui rischi per il Piano di Progetto (PdP)
    \item Pianificazione tecnica e funzionale del Proof of Concept (PoC)
    \item Rotazione dei ruoli per il quinto periodo
\end{itemize}

\newpage
\section{Svolgimento}
\rowcolors{2}{white}{white}
\begin{longtable}{p{5cm}|p{10cm}}
\textbf{Stato avanzamento Analisi dei Requisiti} & Il gruppo espone le difficoltà nella gestione dei casi d'uso su Google Drive e comunica il passaggio a LaTeX per una migliore gestione della numerazione e dei riferimenti. Viene anche discussa la traduzione grafica in diagrammi. \\\\

\textbf{Discussione sui rischi per il Piano di Progetto (PdP)} & Vengono identificati rischi specifici per il PDP, ovvero rischi: tecnologici, interpersonali, organizzativi, requisiti/stakeholder e costi/tempi. \\\\

\textbf{Pianificazione tecnica e funzionale del Proof of Concept (PoC)} & È stata definita la creazione di una repository dedicata denominata "PoC". Le funzionalità primarie includeranno il caricamento di un pacchetto e la visualizzazione dei contenuti. \\\\

\textbf{Rotazione dei ruoli per il quinto periodo} & Definiti i ruoli per il quinto sprint: Sara Gioia Fichera (Responsabile), Bianca Zaghetto (Amministratore), Andrea Masiero e Francesco Balestro (Analisti), Nenad Radulovic (Programmatore), Filippo Compagno e Mattia Oliva Medin (Verificatori). \\
\end{longtable}

\section{Decisioni}
\rowcolors{2}{lightgray}{white}
\begin{tabularx}{\textwidth}{|l|X|}
    \hline
    \textbf{Identificativo} & \textbf{Descrizione} \\
    \hline
    VI\_20251223.d1 & Adozione di LaTeX per la documentazione tecnica dell'Analisi dei Requisiti. \\
    \hline
    VI\_20251223.d2 & Esclusione della funzionalità di ricerca nella prima versione del PoC. \\
    \hline
    VI\_20251223.d3 & Creazione della repository GitHub "PoC" entro l'inizio del prossimo sprint. \\
    \hline
\end{tabularx}

\section{Attività}
\rowcolors{2}{lightgray}{white}
\begin{tabularx}{\textwidth}{|l|X|l|}
    \hline
    \textbf{Incaricato} & \textbf{Task Todo} & \textbf{Identificativo} \\
    \hline
    \textbf{Andrea Masiero} & Stesura Verbale Interno & VI\_20251223.a1 \\
    \hline
    \textbf{Bianca Zaghetto} & Creazione repo PoC e stesura rischi nel PDP & VI\_20251223.a2 \\
    \hline
    \textbf{Nenad Radulovic} & Sviluppo bozze iniziali e architettura PoC & VI\_20251223.a3 \\
    \hline
    \textbf{Andrea Masiero, Francesco Balestro} & Completamento stesura testuale casi d'uso in LaTeX & VI\_20251223.a4 \\
    \hline
    \textbf{Mattia Oliva Medin} & Stesura Verbale Esterno & VI\_20251223.a5 \\
    \hline
\end{tabularx}

\end{document}